\chapter{АНАЛИЗ ПРЕДМЕТНОЙ ОБЛАСТИ}

% -------------------------------------------------------
\section{Когнитивные предубеждения: теоретические основания}
% -------------------------------------------------------

\subsection{Программа эвристик и предубеждений Канемана и Тверски}

Систематическое изучение когнитивных предубеждений как научной проблемы восходит к серии работ Амоса Тверски и Даниэля Канемана, опубликованных в 1970--1980-х годах~\cite{tversky1981framing, tversky1983extensional}. Эта исследовательская программа сформировалась как ответ на доминирующую в то время концепцию \textit{Homo economicus}, предполагающую, что агенты принимают решения на основе строгого математического расчета и максимизации ожидаемой полезности. Опираясь на идею «ограниченной рациональности» (bounded rationality) Герберта Саймона~\cite{simon1955behavioral}, Канеман и Тверски доказали, что в условиях неопределенности и ограниченных вычислительных ресурсов люди систематически нарушают аксиомы рационального поведения. 

Вместо точного байесовского обновления убеждений или полного перебора вариантов человек использует \textit{эвристики}~--- когнитивные «ярлыки» или упрощенные ментальные стратегии. Эволюционно эвристики оправданы: они обеспечивают высокую скорость реакции и когнитивную экономию. Однако в строго формализованных задачах они порождают предсказуемые и воспроизводимые ошибки суждения, называемые когнитивными предубеждениями.

Канеман и Тверски описали три основные эвристики, лежащие в основе широкого класса искажений~\cite{kahneman2011thinking}:
\begin{enumerate}
    \item \textit{Эвристика репрезентативности} состоит в оценке вероятности события по степени его соответствия типичному прототипу или нарративному шаблону без учёта базовых вероятностей (base rates).
    \item \textit{Эвристика доступности} связана с оценкой частоты или вероятности события через лёгкость, с которой соответствующие примеры извлекаются из памяти.
    \item \textit{Эвристика якорения и корректировки} предполагает оценку числовых величин через начальную привязку к некоторому значению (якорю) с последующей, как правило, недостаточной корректировкой.
\end{enumerate}

Для настоящего исследования наиболее релевантна эвристика репрезентативности, поскольку именно она лежит в основе ошибки конъюнкции~--- центральной парадигмы первого эксперимента. Согласно этой эвристике, конъюнктивный вариант ответа кажется более вероятным, если он семантически и нарративно согласуется с описанием объекта, даже если формальные законы теории вероятностей однозначно исключают такой вывод.

Теоретическим обобщением эвристической программы стала двухсистемная модель мышления (Dual-Process Theory)~\cite{evans2003in, kahneman2011thinking}: 
\begin{itemize}
    \item \textbf{Система 1} работает быстро, автоматически, параллельно и ассоциативно. Она опирается на распознавание паттернов и не требует осознанных когнитивных усилий.
    \item \textbf{Система 2} функционирует медленно, последовательно, требует концентрации внимания и отвечает за логический вывод и применение формальных правил.
\end{itemize}

Большинство когнитивных предубеждений возникает из-за эффекта «когнитивной скупости» (cognitive miserliness): Система 1 мгновенно генерирует правдоподобный интуитивный ответ, а Система 2 принимает его без достаточной аналитической проверки. Данная теоретическая рамка имеет прямое значение для интерпретации поведения больших языковых моделей (LLM). Как будет показано далее, авторегрессионная генерация токенов по своей математической природе (сэмплирование на основе априорных распределений) функционально аналогична ряду свойств Системы~1 --- в частности, опоре на априорные распределения и отсутствию явного механизма пересмотра убеждений в процессе единственного прямого прохода. Вместе с тем данная аналогия носит операциональный характер: в отличие от биологической Системы~1, авторегрессионный декодер является последовательным детерминированным процессом, лишённым эволюционно закреплённых эвристик и аффективных компонентов. Техники промптинга, такие как пошаговое рассуждение (Chain-of-Thought), представляют собой попытку искусственно эмулировать отдельные свойства Системы~2 за счёт декомпозиции вывода на явные промежуточные шаги~\cite{wei2022chain}, однако принципиально не изменяют математику декодирования --- веса модели и механизм сэмплирования остаются идентичными.

\subsection{Типология когнитивных предубеждений: вероятностные, логические, индуктивные}

Для целей настоящего исследования когнитивные предубеждения классифицируются по типу нарушаемой нормы рассуждения. Эта классификация позволяет выстроить три экспериментальные парадигмы не как набор изолированных тестов, а как единый прогрессирующий \textit{спектр когнитивной сложности}. Переход от вероятностного суждения к индуктивному открытию правил сопровождается экспоненциальным ростом требований к вычислительным ресурсам агента.

\textbf{1. Вероятностные предубеждения (Одношаговое суждение).}
Возникают при нарушении аксиом теории вероятностей. В данных задачах агенту предоставляется полная информация, и от него требуется лишь оценить статические вероятности. Наиболее известный пример~--- \textit{ошибка конъюнкции} (Conjunction Fallacy)~\cite{tversky1983extensional}. Когнитивная нагрузка здесь минимальна: ошибка возникает исключительно из-за подмены математической вероятности метрикой семантического сходства (нарративной правдоподобности).

\textbf{2. Логические предубеждения (Дедуктивный анализ ограничений).}
Проявляются при нарушении норм формальной логики. Задача выбора карточек Уэйсона (Wason Selection Task)~\cite{wason1966reasoning} является классическим инструментом их диагностики. Дедуктивное рассуждение когнитивно сложнее вероятностного, поскольку требует применения абстрактных синтаксических правил (например, modus tollens: $P \rightarrow Q, \lnot Q \Rightarrow \lnot P$) вопреки семантическому содержанию задачи. Предубеждение подтверждения (Confirmation Bias) и предубеждение соответствия (Matching Bias)~\cite{evans1998matching} заставляют агента выбирать элементы, явно названные в правиле, вместо поиска фальсифицирующих условий.

\textbf{3. Индуктивные предубеждения (Многоходовое обновление убеждений).}
Возникают при нарушении нормативных стандартов формирования и пересмотра гипотез в условиях неполной информации. Задача открытия правила 2-4-6 Уэйсона~\cite{wason1960failure} представляет собой вершину когнитивной сложности в данном спектре. В отличие от дедукции, индукция требует от агента:
\begin{enumerate}
    \item Самостоятельно сформировать начальное пространство гипотез (hypothesis space).
    \item Итеративно генерировать тестовые примеры (сэмплинг).
    \item Осуществлять байесовское обновление убеждений на основе получаемой обратной связи~\cite{tenenbaum1999bayesian}.
\end{enumerate}
Здесь предубеждение подтверждения проявляется в форме «стратегии позитивного тестирования» (positive test strategy) --- склонности тестировать примеры, которые согласуются с текущей гипотезой, игнорируя попытки её фальсифицировать.

Данная типология образует прогрессирующий спектр от задачи пассивного распознавания (вероятностное суждение) к активному исследованию среды (индукция). Это закладывает теоретический фундамент для проверки гипотезы о том, что устойчивость современных ИИ-систем к когнитивным искажениям будет неравномерной и снизится при переходе к задачам, требующим многоходового управления пространством гипотез.

\subsection{Нормативные стандарты рассуждения и отклонения от них}

Для строгой алгоритмической оценки ответов языковой модели необходимо формализовать понятие ошибки. Под нормативным стандартом в настоящей работе понимается математический или логический критерий, относительно которого ответ (или стратегия) агента однозначно классифицируется как правильный или ошибочный. 

\textbf{Стандарт вероятностного суждения (Задача Линды).} 
Нормативный ответ определяется аксиоматикой Колмогорова. Вероятность наступления двух совместных событий не может превышать вероятность наступления любого из них в отдельности: $P(A \cap B) \leq P(A)$ и $P(A \cap B) \leq P(B)$ для любых событий $A$ и $B$. Выбор конъюнктивного варианта ответа (А и Б) при наличии одинарной опции (А) является строгим нарушением экстенсиональной логики классов, вне зависимости от заданного текстового контекста.

\textbf{Стандарт дедуктивного вывода (Задача Уэйсона).} 
Нормативный ответ определяется правилом материальной импликации в логике высказываний. Для проверки истинности условного утверждения формы $P \rightarrow Q$ («Если $P$, то $Q$») необходимо и достаточно проверить те случаи, которые способны сделать утверждение ложным. Таковыми являются только ситуации, где выполняется антецедент ($P$), но не выполняется консеквент ($\lnot Q$). Таким образом, нормативным ответом является выбор карточек $P$ и $\lnot Q$~\cite{wason1966reasoning}. Выбор карточки $Q$ является ошибкой утверждения консеквента, а карточки $\lnot P$~--- ошибкой отрицания антецедента.

\textbf{Стандарт индуктивного вывода (Задача 2-4-6).} 
В отличие от предыдущих парадигм, нормативный стандарт в индуктивной задаче оценивает не одномоментный выбор, а \textit{стратегию тестирования и качество финальной гипотезы}. 
Согласно принципам байесовского оптимального экспериментального дизайна (Bayesian Optimal Experimental Design) и попперовской фальсифицируемости~\cite{oaksford1994rational}, нормативный агент должен максимизировать ожидаемую информационную ценность каждого теста. Это требует генерации таких троек чисел, которые могут \textit{опровергнуть} текущую гипотезу (например, генерация тройки, не соответствующей гипотезе агента, в расчете на получение ответа True). Критерием итогового успеха служит не текстовое совпадение гипотез, а их экстенсиональная эквивалентность, которая в данной работе измеряется через метрику пересечения по объединению (Intersection over Union, IOU) над репрезентативной выборкой числовых троек~\cite{banatt2024wilt}.

Для интерпретации результатов всех трёх экспериментов необходимо обозначить точки отсчёта. \textit{Нормативный агент}~--- идеальный байесовский агент, не допускающий ошибок конъюнкции, выбирающий карточки $P$ и $\lnot Q$ в задаче Уэйсона и генерирующий исключительно фальсифицирующие тройки в задаче 2-4-6 (IOU~$= 1{,}0$, CTR~$= 0$)~--- задаёт верхнюю границу нормативного рассуждения. \textit{Человеческая производительность} фиксирует эмпирически наблюдаемый уровень: $\approx$80\% частота ошибки конъюнкции~\cite{tversky1983extensional}, $\approx$10\% правильных ответов в абстрактной версии задачи Уэйсона~\cite{wason1966reasoning} и преобладание позитивного тестирования в задаче 2-4-6~\cite{wason1960failure}. \textit{Случайный агент}, равновероятно выбирающий из доступных опций, задаёт нижнюю границу осмысленного рассуждения. Результаты тестируемых LLM интерпретируются относительно этих трёх ориентиров.

% -------------------------------------------------------
\section{Анализ проблемы: когнитивные предубеждения в LLM}
% -------------------------------------------------------

\subsection{Природа предубеждений в языковых моделях: memorization vs. emergent bias}

Принципиальный методологический вопрос при тестировании больших языковых моделей на когнитивных задачах состоит в интерпретации наблюдаемых ошибок. Существуют два конкурирующих объяснения того, почему искусственная нейронная сеть может демонстрировать поведенческие паттерны, характерные для человеческих когнитивных искажений.

Согласно гипотезе \textit{memorization} (запоминания), модель воспроизводит правильный или ошибочный ответ исключительно потому, что видела данную конкретную задачу или её близкие аналоги в обучающем корпусе. Классические парадигмы, такие как задача Линды, задача Уэйсона и задача 2-4-6, широко представлены в психологической литературе, на которой обучаются современные LLM. Таким образом, успешное решение канонической задачи может не свидетельствовать о наличии абстрактного логического рассуждения, равно как и воспроизведение ошибки может отражать лишь статистическое запоминание примеров человеческой нерациональности из обучающих данных~\cite{zhou2023benchmark}.

Согласно гипотезе \textit{emergent bias} (эмерджентного предубеждения), искажения возникают как неотъемлемое следствие самой природы авторегрессионного обучения на массивах человеческих текстов. Модель усваивает не только изолированные факты, но и латентные структуры человеческого рассуждения, включая систематические ассоциативные смещения. В этом случае когнитивные предубеждения должны устойчиво воспроизводиться на новых, контаминационно чистых (out-of-distribution) стимулах, сохраняющих логическую структуру классических задач.

Эмпирические данные остаются противоречивыми. Ряд исследований демонстрирует свидетельства в пользу memorization: модели успешно решают канонические версии задач, однако их точность резко падает при изменении поверхностных признаков стимула~\cite{zhou2023benchmark, macmillan2024irrationality}. Dasgupta и соавторы~\cite{dasgupta2024language}, напротив, выявили контент-эффекты на специально сконструированных деконтаминированных стимулах, что указывает на emergent bias. Настоящее исследование фокусируется на измерении итогового поведения модели в контролируемых условиях, рассматривая предубеждения как системную уязвимость вне зависимости от первичного механизма их усвоения. Данный методологический выбор является сознательным ограничением: применяемый экспериментальный дизайн позволяет зафиксировать факт наличия и интенсивность предубеждения, однако не позволяет атрибутировать его источник --- memorization или emergent bias. Разграничение этих двух механизмов требует отдельного контролируемого исследования с использованием верифицированно деконтаминированных корпусов и составляет перспективу дальнейших работ.

\subsection{Архитектурные и математические предпосылки эвристик в трансформерах}

Несмотря на поведенческое сходство ошибок LLM с искажениями человека, для понимания их природы необходимо обратиться к математическим механизмам архитектуры Transformer. Исследования в области механистической интерпретируемости (mechanistic interpretability) за 2023--2025 гг. предоставляют сходящиеся доказательства того, что авторегрессионное моделирование по умолчанию действует как «быстрый эвристический предиктор», структурно эквивалентный Системе 1~\cite{zhang2025identifying}. 

Когнитивные предубеждения в LLM обусловлены двумя взаимодействующими математическими структурами:

\textbf{1. Доминирование априорных распределений в пространстве логитов (Эвристика доступности).} 
В декодерных архитектурах вектор скрытого состояния проецируется в пространство логитов через матрицу деэмбеддинга (unembedding matrix), после чего применяется функция softmax. Stolfo и соавторы~\cite{stolfo2024confidence} показали существование специфических «нейронов частотности» (token frequency neurons), которые формируют аддитивные сдвиги логитов, выровненные по направлению логарифма базовой частоты токенов ($v_{\text{freq}} \approx \log p_{\text{freq}} - \text{mean}(\log p_{\text{freq}})$). Поскольку softmax применяется к логитам, такое аддитивное смещение математически эквивалентно мультипликативному перевзвешиванию вероятностей на основе униграммного распределения. Это означает, что в условиях неопределённости модель архитектурно склонна возвращаться к базовым (наиболее частотным) ответам, игнорируя контекстную логику задачи. Следует оговориться, что данный результат получен преимущественно на моделях семейства GPT-2/GPT-3 и требует отдельной верификации применительно к современным RLHF-выровненным архитектурам большего масштаба. Тем не менее общий принцип частотного смещения логитов подтверждается независимыми работами в области механистической интерпретируемости~\cite{zhang2025identifying}. Данный механизм является функциональным аналогом эвристики доступности.

\textbf{2. Softmax Attention как ядерная регрессия (Эвристика репрезентативности).} 
Анализ формы softmax-внимания показывает, что оно математически эквивалентно экспоненциальной ядерной сглаживающей регрессии (ядерной оценке Надарая-Уотсона)~\cite{collins2024softmax, han2023explaining}. В контексте in-context learning скалярные произведения (dot-products) ключей и запросов выполняют роль функции сходства. В результате механизм внимания по своей природе отдает приоритет токенам с высоким семантическим сходством с текущим контекстом. Это прямое воплощение эвристики репрезентативности: модель оценивает вероятность следующего токена по степени его семантической близости к нарративу (например, связывая «феминизм» с описанием Линды), а не по строгим правилам конъюнкции. Следует оговориться, что данная аналогия строго выводится лишь для упрощённого случая одноголового внимания на одномерных последовательностях; в полноценных многоголовых трансформерах соответствие является приближённым и функциональным, а не математически точным~\cite{collins2024softmax}.

Кроме того, исследования логических цепей в LLM~\cite{fu2024implies} демонстрируют, что структурные логические операторы («И», «ИЛИ») обладают низкой салиентностью (low-salience) для стандартных голов внимания. Если в модели не активируются специализированные головы внимания к операторам, она переходит в режим «ленивого рассуждения» (lazy reasoning), опираясь исключительно на семантическое сходство токенов-существительных, что неизбежно ведет к логическим ошибкам.

\subsection{Влияние методов выравнивания (RLHF) и склонность к угождению пользователю}

Помимо базовой архитектуры, существенным источником индуктивных и логических искажений (в частности, предубеждения подтверждения) является этап пост-тренировочного выравнивания (Post-Training Alignment). 

Современные LLM проходят обучение с подкреплением на основе отзывов людей (RLHF) или прямую оптимизацию предпочтений (DPO). Функция вознаграждения в этих методах максимизирует сходство ответов модели с предпочтениями асессоров. Однако люди-асессоры сами подвержены когнитивным искажениям. Оптимизируясь под человеческое одобрение, модели приобретают свойство \textit{угождать пользователю} --- склонности соглашаться с утверждениями пользователя, даже если они логически ошибочны, и избегать конфронтационных или фальсифицирующих ответов~\cite{batista2026rational}. 

Это имеет критические последствия для задач индуктивного вывода (таких как парадигма 2-4-6). Вместо активного поиска фальсифицирующих примеров, необходимых для нормативного байесовского рассуждения, выровненные (aligned) модели склонны генерировать подтверждающие примеры, удовлетворяя имплицитную гипотезу пользователя. Таким образом, RLHF может вносить дополнительный вклад в усиление Confirmation Bias. Вместе с тем разграничение этого эффекта от предубеждений, унаследованных на этапе предобучения (pretraining), остаётся открытым исследовательским вопросом: имеющиеся эмпирические данные носят корреляционный характер и не позволяют установить направление причинно-следственной связи~\cite{batista2026rational}. Разрешение данного вопроса требует контролируемого сравнения base-моделей и их RLHF-выровненных аналогов на идентичных стимульных наборах.

\subsection{Проблема контаминации обучающих данных}

Контаминация обучающих данных (training data contamination) --- одна из центральных методологических угроз при оценке когнитивных способностей LLM. В отличие от стандартных бенчмарков (например, MMLU), где контаминация всегда ведет к завышению точности, в когнитивных задачах она вызывает разнонаправленные эффекты. Модель может запомнить как правильный нормативный ответ (завышение метрик), так и человекоподобный ошибочный паттерн из учебников по психологии (искусственное возникновение bias).

Для обхода этой проблемы предлагались различные стратегии: изменение нарративного фрейминга~\cite{suri2024large}, динамическая генерация синтетических задач~\cite{jiang2024peek} или хеширование смысловых токенов~\cite{chadimova2024meaningless}. В настоящем исследовании контаминация строго контролируется на уровне каждого эксперимента: в задаче Линды применяется факториальный дизайн с подстановкой оригинальных биографий и демографических токенов; в задаче Уэйсона введены условия «фантастических социальных контрактов»; для задачи 2-4-6 сгенерирован независимый верифицированный датасет правил, исключающий прямое пересечение с корпусами.

\subsection{Ограничения существующих подходов и обоснование метрик}

Анализ актуальной литературы выявляет несколько системных ограничений в существующей методологии тестирования LLM на когнитивные предубеждения.

\textbf{Узость охвата моделей.} Большинство работ тестируют одну-две проприетарные модели (как правило, семейства GPT), что не позволяет отделить архитектурные артефакты от фундаментальных свойств трансформеров. Настоящее исследование расширяет охват до девяти моделей различных семейств и размеров.

\textbf{Иллюзия бинарных метрик (Accuracy).} Преобладающим подходом в бенчмарках (например, BIG-bench, CogBench) является подсчёт доли правильных ответов (Accuracy или Exact Match). Данный подход фундаментально ограничен для когнитивных исследований. Бинарная метрика не способна различить модель, угадавшую правильный ответ с уверенностью 51\%, и модель, логически выведшую его с уверенностью 99\%. Более того, уверенность модели в \textit{ошибочном} ответе является важнейшим индикатором силы предубеждения. В связи с этим в настоящей работе вводится непрерывная метрика силы предубеждения (\textit{Severity of Bias}), основанная на калибровке уверенности, а также метрика экстенсиональной эквивалентности (IOU).

\textbf{Отсутствие оценки стабильности.} Большинство классических тестов являются однократными (single-shot). Однако LLM обладают высокой стохастичностью. В данной работе стабильность ответов измеряется как самостоятельная характеристика через систематические перестановки (permutations) позиций стимулов, что позволяет отделить робастное логическое рассуждение от случайного угадывания.

\subsection{Практическая значимость и целевая аудитория}

Результаты настоящего исследования ориентированы на три группы потребителей. 
Во-первых, исследователи в области alignment и безопасности ИИ получают 
воспроизводимый фреймворк для профилирования когнитивных уязвимостей конкретных 
моделей перед их развёртыванием в критических приложениях. Во-вторых, инженеры 
и архитекторы LLM-систем могут использовать предложенные метрики (Severity of 
Bias, AUC-IOU) для сравнительного отбора моделей в задачах, требующих надёжного 
вероятностного и логического вывода~--- например, в медицинской диагностике, 
юридическом анализе или финансовом консультировании. В-третьих, разработчики 
бенчмарков получают методологический шаблон для построения многопарадигмальных 
оценок когнитивных способностей ИИ, выходящих за рамки бинарных метрик точности.

Конкретным сценарием применения является следующий: перед интеграцией LLM в 
систему поддержки принятия решений инженерная команда запускает предложенный 
пайплайн на целевой модели и получает количественный профиль её когнитивных 
уязвимостей в разрезе трёх типов рассуждения. Это позволяет либо выбрать 
альтернативную модель с меньшим bias-профилем, либо целенаправленно применить 
митигирующие техники (например, CoT-промптинг) именно для тех парадигм, где 
модель наиболее уязвима.

% -------------------------------------------------------
\section{Обзор смежных исследований и существующих бенчмарков}
% -------------------------------------------------------

\subsection{Ошибка конъюнкции, эвристика репрезентативности и демографический токен-bias}

Задача Линды была одной из первых классических психологических парадигм, применённых к тестированию языковых моделей. Suri и соавторы~\cite{suri2024large} показали, что GPT-3.5 воспроизводит ошибку конъюнкции с частотой $\approx$80\%, сопоставимой с человеческой. Macmillan-Scott и Musolesi~\cite{macmillan2024irrationality} подтвердили аналогичную картину на нескольких моделях в рамках масштабного бенчмарка когнитивных искажений. 

Однако более глубокий механистический подход к этой проблеме требует анализа \textit{токен-bias} --- предрасположенности модели выбирать ответ с более высокой априорной вероятностью в языковом корпусе вне зависимости от логической структуры задачи. В этом контексте особую важность приобретают исследования 2024--2025 годов, изучающие влияние демографических маркеров на вероятностный и логический вывод.

Shafiei и соавторы (2025)~\cite{shafiei2025more} в бенчмарке «More or Less Wrong» продемонстрировали, что замена нейтрального субъекта на демографически нагруженный токен (например, «женщина», «афроамериканец») в строго математических задачах вызывает резкий дрейф предсказаний (directional error drift) в сторону стереотипных ассоциаций. Модель Claude 3.7 Sonnet изменяла частоту ошибок на 27 процентных пунктов исключительно из-за замены демографического токена.

Jiang и соавторы (EMNLP 2024)~\cite{jiang2024peek} применили методологию контрфактических замен (string replacement) к логическим парадоксам. Заменяя известные имена на общие существительные при сохранении логической структуры конъюнкции, авторы использовали статистический тест МакНемара (McNemar test) и доказали, что успешность LLM часто является артефактом токен-bias, а не следствием истинного рассуждения (genuine reasoning). 

An и соавторы (PNAS 2025)~\cite{an2025measuring} подтвердили этот эффект на задачах скоринга резюме, где замена исключительно имени кандидата (сигнализирующего о расе и гендере) статистически значимо сдвигала оценку LLM (например, у GPT-3.5 Turbo), несмотря на жесткую фиксацию объективных параметров. Механистически этот эффект объясняется работами Yang et al. (2025)~\cite{yang2025bias}, показавшими, что специфические головы внимания (attention heads) активируются на демографические токены и искажают логический вывод в пользу стереотипного эмбеддинга.

Настоящее исследование (Эксперимент 1) напрямую развивает эту линию, объединяя проверку ошибки конъюнкции с факториальным дизайном демографических подстановок имен (token-bias) в оригинальных биографических виньетках.

\subsection{Дедуктивное рассуждение и задача выбора карточек Уэйсона}

Задача Уэйсона является наиболее изученной логической парадигмой в исследованиях LLM. В то время как абстрактная версия (A, B, 4, 7) часто решается современными моделями успешно, исследователи сходятся во мнении, что это результат запоминания (memorization). Macmillan-Scott и Musolesi~\cite{macmillan2024irrationality} показали, что ответы LLM на задачу Уэйсона существенно нестабильны: незначительные изменения формулировки или порядка предъявления карточек ведут к резким изменениям точности.

Dasgupta и соавторы~\cite{dasgupta2024language} провели методологически строгое исследование деконтаминированных стимулов и обнаружили, что крупные модели воспроизводят контент-эффекты, аналогичные человеческим. В частности, социальные контракты (деонтическая логика «разрешений и обязательств») решаются точнее, чем абстрактные правила, что согласуется с теорией прагматических схем рассуждения Ченга и Холиоака~\cite{cheng1985pragmatic}.

Масштабный бенчмарк Macmillan-Scott и Musolesi~\cite{macmillan2024irrationality} выявил существенную межмодельную вариабельность: одни LLM демонстрируют точность, близкую к потолку, другие подвержены систематическому предубеждению соответствия (matching bias), описанному Эвансом~\cite{evans1998matching}. Важным выводом их работы стала фиксация высокой нестабильности ответов при повторных запросах, что ставит под сомнение надежность точечных (single-shot) оценок логических способностей ИИ. 

Для преодоления этих ограничений в Эксперименте 2 настоящего исследования вводится уникальное условие «фантастических социальных контрактов» (позволяющее отделить распознавание деонтической структуры от эффекта знакомости контента), а также измеряется стабильность ответов на основе систематических перестановок.

\subsection{Индуктивное открытие правил и задача 2-4-6}

Исследования индуктивного рассуждения LLM развиты значительно слабее по сравнению с вероятностными и дедуктивными задачами. Индукция требует от агента генерации гипотез, их активной проверки и обновления убеждений, что является когнитивно сложным процессом (System 2).

Banatt и соавторы~\cite{banatt2024wilt} сделали важный шаг, разработав бенчмарк WILT --- первый систематический инструмент оценки LLM в многоходовой задаче открытия правил. Их результаты показали, что хотя более крупные модели успешнее открывают правила, абсолютно все модели проявляют признаки предубеждения подтверждения (Confirmation Bias): доля подтверждающих тестов (positive testing) катастрофически превышает нормативную байесовскую базу. 

Batista и Griffiths~\cite{batista2026rational} связали это поведение с алгоритмическим выравниванием (RLHF). Оптимизируясь на одобрение, модели неосознанно воспроизводят предпочтение подтверждающей информации как следствие смещения функции вознаграждения. Подобные выводы подтверждаются бенчмарком CogBench~\cite{coda2024cogbench}, который зафиксировал существенный разрыв между LLM и нормативными агентами именно в задачах, требующих активного поиска фальсифицирующих свидетельств.

Настоящее исследование (Эксперимент 3) расширяет эту парадигму, внедряя строгую метрику экстенсиональной эквивалентности гипотез (Intersection over Union, IOU) на базе симуляции Монте-Карло, что позволяет математически точно оценивать траекторию индуктивного поиска, абстрагируясь от синтаксических различий в текстовых формулировках правил моделью.

\subsection{Сравнительные бенчмарки и формирование исследовательского ландшафта}

Несмотря на рост числа исследований когнитивных способностей ИИ, большинство работ сфокусировано на одной парадигме или узком наборе моделей. Ряд исследователей пытались создать интегрированные бенчмарки. Koo и соавторы~\cite{koo2024benchmarking} разработали систему оценки когнитивных предубеждений для LLM в роли судей (LLM-as-a-judge). Проекты вроде CogBench~\cite{coda2024cogbench} агрегируют задачи из теории игр, обучения с подкреплением и вероятностного вывода. 

Однако общим недостатком существующих агрегированных бенчмарков (включая традиционные MMLU и BIG-bench) является изоляция типов рассуждения друг от друга и опора на жесткие бинарные метрики точности. Настоящее исследование заполняет этот пробел, рассматривая три классические парадигмы (вероятность, дедукцию, индукцию) как единый концептуальный спектр. Применяя к ним оригинальные непрерывные метрики (Severity of Bias, AUC-IOU) на наборе из девяти различных архитектур, данная работа переходит от фиксации единичных ошибок к системному профилированию когнитивной уязвимости современных LLM.

% -------------------------------------------------------
\section{Обоснование исследовательских гипотез}
% -------------------------------------------------------

На основании проведенного теоретического анализа архитектурных особенностей LLM и обзора смежных исследований, в настоящей работе формулируются четыре исследовательские гипотезы. Они охватывают аспекты универсальности когнитивных искажений, влияния масштаба моделей, когнитивной сложности решаемых задач и алгоритмической справедливости (токен-bias).

\subsection{Гипотеза об универсальности предубеждений}

\textbf{Гипотеза H1 (Универсальность).} \textit{Все тестируемые языковые модели демонстрируют статистически значимое предубеждение подтверждения (Confirmation Bias) и подверженность эвристике репрезентативности во всех трёх парадигмах относительно формального нормативного стандарта.}

\textbf{Обоснование:} Данная гипотеза опирается на два фундаментальных механизма, описанных в Разделе 1.2. Во-первых, архитектура Transformer математически предрасположена к Системе 1: механизм softmax-внимания действует как ядерная сглаживающая регрессия~\cite{han2023explaining}, отдавая приоритет семантически близким токенам (эвристика репрезентативности) в ущерб структурным логическим операторам. Во-вторых, обучающие корпуса состоят из человеческих текстов, изобилующих нерациональными паттернами. Процедура выравнивания (RLHF) дополнительно пенализирует конфронтационное поведение модели, искусственно поощряя ее угождать пользователю~\cite{batista2026rational} --- стремление подтверждать имплицитные предположения пользователя, что неизбежно ведет к предубеждению подтверждения в логических и индуктивных задачах.

\subsection{Гипотеза о масштабировании и внутрисемейном парадоксе}

\textbf{Гипотеза H2 (Масштабирование).} \textit{В среднем более крупные модели демонстрируют более низкий уровень когнитивных предубеждений за счет лучших способностей к эмуляции Системы 2. Однако данная зависимость не является строго монотонной: внутри модельных семейств может наблюдаться парадокс, при котором меньшая модель демонстрирует меньший bias, чем крупная.}

\textbf{Обоснование:} Традиционно в ML предполагается, что увеличение количества параметров ведет к монотонному улучшению метрик. Однако в задачах, подверженных когнитивным искажениям, часто наблюдается эффект «U-образного масштабирования» (U-shaped scaling) или обратного масштабирования (inverse scaling)~\cite{dasgupta2024language}. С одной стороны, масштаб улучшает базовые навыки рассуждения (Chain-of-Thought). С другой стороны, более крупная модель обладает более богатым и детализированным пространством эмбеддингов, что позволяет ей лучше улавливать поверхностные стереотипы и нарративную связность текста. Таким образом, крупная модель может успешнее распознать «ловушку» эвристики репрезентативности и... увереннее в неё попасться, опираясь на свои глубокие априорные знания, в то время как малая модель ошибется случайно или выберет базовую вероятность.

\subsection{Гипотеза о спектральной природе предубеждений}

\textbf{Гипотеза H3 (Спектральность).} \textit{Интенсивность когнитивных предубеждений (отклонений от нормативного стандарта) систематически возрастает от вероятностной парадигмы к логической и достигает максимума в индуктивной, отражая экспоненциальное возрастание когнитивной сложности требуемого рассуждения.}

\textbf{Обоснование:} Как было показано в Разделе 1.1.2, три выбранные парадигмы не равнозначны. 
\begin{itemize}
    \item Одношаговое вероятностное суждение (задача Линды) требует лишь подавления априорного токен-bias.
    \item Дедуктивный анализ (задача Уэйсона) требует применения абстрактного правила (modus tollens) поверх семантического контекста.
    \item Индуктивное открытие правил (задача 2-4-6) требует многоходового удержания контекста, генерации пространства гипотез (hypothesis space) и активного фальсифицирующего сэмплинга.
\end{itemize}
Поскольку авторегрессионные модели испытывают фундаментальные трудности с долгосрочным планированием и активным поиском опровержений (отсутствие истинного внутреннего цикла обратной связи), ожидается, что метрики фальсификации (например, Confirming Test Rate) в задаче 2-4-6 будут стабильно демонстрировать максимальный уровень предубеждения подтверждения у всех моделей.

\subsection{Гипотеза о демографическом токен-bias}

\textbf{Гипотеза H4 (Демографический bias).} \textit{Демографические переменные, кодируемые именем персонажа, оказывают статистически значимое влияние на интенсивность вероятностной ошибки конъюнкции, при этом различия предположительно будут более выражены по оси расы, чем по осям возраста и пола. Данное утверждение само является проверяемой гипотезой и подлежит эмпирической верификации по результатам эксперимента --- направление эффекта выводится из теоретических соображений о топологии эмбеддингов, но не может считаться установленным до получения экспериментальных данных.}

\textbf{Обоснование:} На основе методологии контрфактических замен (string replacement), успешно примененной в работах Jiang et al. (2024)~\cite{jiang2024peek} и Shafiei et al. (2025)~\cite{shafiei2025more}, предполагается, что имена собственные функционируют в LLM не как нейтральные идентификаторы, а как плотные семантические векторы. Имена, статистически ассоциированные с определенными расовыми или этническими группами в обучающем корпусе, обладают специфической топологией в пространстве эмбеддингов. 
Внимание (Attention heads) модели, действуя как механизм поиска по сходству, может активировать различные стереотипные ассоциации в зависимости от подаваемого имени, искажая чистоту математического вывода. Возраст, задаваемый числовым значением, обладает меньшей социокультурной «гравитацией» в пространстве эмбеддингов, поэтому теоретически ожидается, что расовые маркеры (имена) спровоцируют более сильные флуктуации в метриках bias. Вместе с тем прямое количественное сравнение силы расового и гендерного эффектов в цитируемых работах~\cite{jiang2024peek, shafiei2025more} не проводилось в едином факториальном дизайне, что не позволяет считать данное направленное предсказание строго выводимым из литературы --- оно формулируется как гипотеза, требующая проверки.

% -------------------------------------------------------
\section{Требования к экспериментальному исследованию}
% -------------------------------------------------------

На основании проведенного критического анализа существующих подходов и выявленных методологических ограничений (Разделы 1.2.4--1.3.4), для эмпирической проверки выдвинутых гипотез формулируется комплекс строгих требований к экспериментальному фреймворку. Эти требования призваны обеспечить научную чистоту, воспроизводимость результатов и статистическую значимость выводов при оценке когнитивных предубеждений в LLM.

\subsection{Требования к датасетам и контролю контаминации}

\textbf{1. Оригинальность стимулов и деконтаминация (Out-of-Distribution testing).} 
Ключевым требованием является минимизация риска \textit{memorization}. Все биографические виньетки для задачи Линды (Эксперимент 1) должны быть сгенерированы вручную с использованием факториального дизайна (factorial design) и не воспроизводить текст оригинальной задачи Тверски и Канемана. Для задачи Уэйсона (Эксперимент 2) и задачи 2-4-6 (Эксперимент 3) стимульные сеты должны включать деконтаминированные условия с оригинальным абстрактным или вымышленным содержанием (например, условие «незнакомых фантастических контрактов»).

\textbf{2. Структурное единообразие и изоляция переменных.} 
Каждая единица датасета должна быть сконструирована по жесткому синтаксическому шаблону. Операции контрфактической замены (string replacement) демографических токенов или логических операторов должны производиться изолированно, исключая появление конфаундеров (confounders), связанных с вариабельностью длины промпта, его лексической сложности или нарративного стиля.

\textbf{3. Объем выборки и статистическая мощность (Statistical Power).} 
Датасеты должны обеспечивать достаточный объем наблюдений для достижения статистической значимости применяемых непараметрических тестов. Для теста МакНемара на согласованных парах необходимо минимум 30 дискордантных пар на каждое сравниваемое условие. Оценка демографических эффектов (H4) требует достаточного числа факториальных подстановок на каждую исследуемую страту (расовую или гендерную группу).

\textbf{4. Верификация консистентности.} 
Каждый элемент генеративного датасета должен пройти автоматизированную или экспертную проверку на отсутствие логических противоречий, непреднамеренных тавтологий и функциональных дубликатов (проверка через экстенсиональную эквивалентность сгенерированных правил в Эксперименте 3).

\subsection{Требования к метрикам и статистическому выводу}

\textbf{1. Переход от бинарной классификации к вероятностным метрикам.} 
Ввиду ограниченности метрики Accuracy (как обосновано в 1.2.4), экспериментальный фреймворк обязан использовать непрерывные метрики оценки. Фреймворк должен фиксировать самооценку уверенности модели (confidence) и агрегировать её в метрику \textit{Severity of Bias}, что позволит обнаруживать латентные предубеждения даже в режиме «потолка бинарной точности» (ceiling effect).

\textbf{2. Оценка экстенсиональной эквивалентности (Intersection over Union).} 
Для оценки индуктивного рассуждения (многоходовой генерации гипотез) текстовое совпадение (exact match) неприменимо из-за синтаксической вариативности LLM. Требуется внедрение метрики IOU на основе Монте-Карло сэмплирования, позволяющей математически строго сравнивать сгенерированные моделью предикаты (лямбда-выражения) с истинным скрытым правилом.

\textbf{3. Парный дизайн экспериментов (Paired Experimental Design).} 
Для изоляции эффекта конкретного токена или семантической связи необходимо использовать внутрисубъектный дизайн (within-subject design). Сравнение условий (например, \textit{correlated} vs. \textit{uncorrelated} опции; прямая vs. обратная перестановка карточек) должно производиться с использованием точного теста МакНемара~\cite{mcnemar1947note}, обладающего высокой мощностью для парных бинарных наблюдений.

\textbf{4. Контроль частоты ложных открытий (FDR Control).} 
При проведении множественных сравнений (например, между 9 моделями или 5 демографическими группами) в обязательном порядке должна применяться процедура коррекции $p$-значений, такая как метод Бенджамини-Хохберга (Benjamini-Hochberg)~\cite{benjamini1995controlling}, для минимизации риска ошибок первого рода.

\subsection{Требования к инфраструктуре и набору тестируемых моделей}

\textbf{1. Архитектурное разнообразие.} 
Для обеспечения высокой генерализуемости (generalizability) выводов и проверки гипотезы о масштабировании (H2), выборка моделей должна охватывать:
\begin{itemize}
    \item Различных провайдеров (OpenAI, Anthropic, xAI, Google, Alibaba).
    \item Различные парадигмы доступа (закрытые коммерческие API и открытые open-weight модели).
    \item Широкий спектр параметров (от компактных моделей класса 7B--12B до масштабных систем 70B+).
\end{itemize}

\textbf{2. Воспроизводимость и управление состоянием (Reproducibility).} 
Экспериментальный пайплайн должен гарантировать детерминированность среды (насколько это позволяет API поставщиков). Требуется фиксация параметра температуры (например, $T=0$ для задач классификации; $T=1.0$ при тестировании робастности к перестановкам). Сырые ответы моделей (raw completions) и логи нарушения форматов вывода должны сохраняться изолированно от расчетных производных метрик для обеспечения аудируемости результатов.

\textbf{3. Единообразие пайплайна (API Abstraction).} 
Для исключения артефактов, вызванных различиями в препроцессинге промптов или системных сообщениях у разных провайдеров, все модели должны запрашиваться через единый унифицированный интерфейс (например, прокси-маршрутизатор OpenRouter) с идентичной конфигурацией токенизации и системных ролей.